\documentclass[11pt,oneside]{scrbook}
\usepackage[a4paper,margin=2.4cm,marginparwidth=3.2cm,marginparsep=0.5cm]{geometry}
\usepackage{../common/rlstudy}
\usepackage{longtable}

\title{{\color{rlblue}\textbf{STUDY-RL}}\\\Huge Deep Reinforcement Learning\\\Large 24 core weeks plus applied audit and causal biology}
\author{Foundations \textperiodcentered{} Deep RL \textperiodcentered{} LLM RL \\ Safety-critical RL \textperiodcentered{} Matching \textperiodcentered{} Causal biology}
\date{\today}

\begin{document}
\maketitle
\tableofcontents
\newpage

\chapter*{Preface}
\addcontentsline{toc}{chapter}{Preface}

This document is the master syllabus for \textsc{STUDY-RL}, a 24-week self-paced curriculum
designed for a learner who is fluent in modern deep learning but new to reinforcement learning.
An additional applied audit in week 25 connects the core material to a synthetic
ATS experiment and an application review. Week 26 adds causal inference and
discovery across biological environments, with task-aware enhancement of noisy,
sparse measurements. It extends the OPE and offline-data lessons rather than
repeating their policy-value estimators.
It threads three goals through every week:
\begin{enumerate}
  \item \textbf{Applied mathematics} --- connect assumptions, derivations, worked examples, and code; distinguish exact results from approximations.
  \item \textbf{Modern research topics} --- selected readings on DPO/GRPO/RLVR, DreamerV3, EfficientZero v2, BBF, CrossQ, diffusion policies, and LLM agents. Topic coverage is not a claim that every method is implemented or remains the best current choice.
  \item \textbf{Two applied threads} — \emph{safety-critical RL} (one callout per lecture plus a dedicated 3-week track) and \emph{matching/ranking} (insurance lead matching and recruitment candidate matching, woven through bandits, OPE, fairness).
\end{enumerate}

Each week includes a teaching guide (\texttt{teaching\_guide.md}), a Tufte-flavored PDF (\texttt{notes.pdf}), a Jupyter
workbook with graded exercises (\texttt{workbook.ipynb}), curated extended reading
(\texttt{extended\_reading.pdf}), and a mini-project with tests (\texttt{project/}, \texttt{tests/}).

\chapter*{How to use this curriculum}
\addcontentsline{toc}{chapter}{How to use}

\begin{itemize}
  \item About \textbf{10 hours per week}: 2--3 h reading the PDF, 3--4 h on the workbook, 3--4 h on the project.
  \item Start with \texttt{STUDY\_GUIDE.md} in the repository root, then the
        weekly teaching guide. Predict an update, calculate it by hand, and
        explain the result before running it.
  \item Attempt workbook exercises before opening hints or reference solutions.
        Some supplied cells call the reference directly: these demonstrate a
        solution but do not test an independent implementation.
  \item Check both execution and understanding: reproduce a numerical example,
        explain the assumptions, test a changed input, and compare a project
        result with a baseline. Passing tests alone does not establish mastery.
\end{itemize}

\chapter{Curriculum overview}

\section{Tracks}
\begin{description}
  \item[T1 — Foundations] (Weeks 1--4) Bandits, MDPs, Bellman equations, dynamic programming, Monte Carlo, TD/SARSA/Q-learning.
  \item[T2 — Deep value-based] (Weeks 5--6) DQN, Double/Dueling/PER, distributional RL, Rainbow, BBF.
  \item[T3 — Policy methods] (Weeks 7--10) REINFORCE, A2C, GAE, TRPO, PPO, SAC/TD3, CrossQ.
  \item[T4 — Bandits \& matching (applied)] (Weeks 11--12) LinUCB/TS, OPE, CLTR, fairness, two-sided markets.
  \item[T5 — Multi-agent RL] (Week 13) MADDPG, MAPPO, self-play, PSRO, AlphaZero case study.
  \item[T6 — Offline RL] (Weeks 14--15) CQL, IQL, AWR, TD3+BC, Decision Transformer, diffusion policies.
  \item[T7 — Model-based RL] (Week 16) Dreamer V3, MuZero/EfficientZero V2, TD-MPC2, DIAMOND.
  \item[T8 — Exploration \& meta-RL] (Week 17) RND, ICM, BYOL-Explore, MAML/PEARL, in-context RL.
  \item[T9 — LLM RL] (Weeks 18--21) RLHF/PPO, DPO/IPO/KTO/ORPO/SimPO, GRPO/RLOO/RLVR, PRMs and reasoning RL.
  \item[T10 — Safety-critical RL] (Weeks 22--24) CMDPs, CPO, shielding, formal verification, scalable oversight, capstone.
  \item[T11 --- Applied audit] (Week 25) Synthetic ATS experiments, evaluation assumptions, and an RL enhancement audit.
  \item[T12 --- Causal biology] (Week 26) Confounding, identification versus discovery, Markov equivalence, multi-environment causal learning, and task-aware enhancement.
\end{description}

\section{Week-by-week schedule}
The topic column includes algorithms studied in the readings. The project
column describes the bundled implementation; each weekly teaching guide
explains its simplifications. A successful toy run is not a reproduction of a
paper's benchmark result. Larger benchmark environments require separate
experiment configuration and evaluation.
\begin{longtable}{@{}p{1.4cm}p{6.4cm}p{6.0cm}@{}}
\toprule
\textbf{Week} & \textbf{Topic} & \textbf{Bundled mini-project} \\
\midrule
\endhead
01 & MDPs, bandits, the RL framing & Multi-armed bandit testbed + UCB/TS \\
02 & Bellman equations, dynamic programming & Value/policy iteration on GridWorld \\
03 & Monte Carlo prediction \& control & Blackjack on-policy MC \\
04 & TD, SARSA, Q-learning, Expected SARSA & CliffWalking + Taxi-v3 \\
05 & Function approximation \& Deep Q-Networks & DQN on LunarLander \\
06 & Rainbow DQN + modern tricks (BBF, NoisyNets) & Rainbow-lite ablations on LunarLander \\
07 & Policy gradients \& REINFORCE & REINFORCE + baselines on CartPole \\
08 & A2C, GAE, natural PG, TRPO & A2C+GAE on LunarLander-Continuous \\
09 & Proximal Policy Optimization (deep dive) & PPO on BipedalWalker + the 37 details checklist \\
10 & SAC, TD3, CrossQ for continuous control & SAC on Pendulum \\
11 & Bandits + contextual bandits + OPE & Off-policy evaluation of a synthetic insurance lead--market matcher \\
12 & RL for ranking / matching (counterfactual LTR, fairness, two-sided) & ATS-style candidate--job ranker with fairness audit \\
13 & Multi-agent RL + AlphaZero case study & Tic-Tac-Toe search/self-play and small MAPPO coordination/spread tasks \\
14 & Offline RL fundamentals (CQL, IQL, AWR, TD3+BC) & Tabular IQL, CQL, and naive Q on synthetic GridWorld logs \\
15 & Decision Transformer + diffusion policies & Small DT on synthetic GridWorld trajectories and numerical diffusion exercises \\
16 & Model-based RL (Dreamer V3, MuZero, TD-MPC2) & Dreamer-lite on CartPole \\
17 & Exploration + meta-RL + in-context RL & RND and a full-information MAML-inspired bandit toy \\
18 & RLHF + PPO for LLMs (from scratch and via TRL) & Character-model SFT, reward learning, and PPO; optional real LoRA lab \\
19 & DPO + IPO + KTO + ORPO + SimPO & Character-model preference objectives; optional real DPO lab \\
20 & GRPO + RLOO + RLVR + process reward models & Group-relative updates on synthetic arithmetic prompts \\
21 & Reward hacking, length bias, KL control, scalable oversight & Reward-hack lab + Constitutional-AI mini-loop \\
22 & Constrained MDPs, Lagrangian methods, PPO-Lagrangian, CPO & PPO-Lagrangian on synthetic LavaGridWorld \\
23 & Safe exploration, shielding, robust RL, formal methods & Model-based grid shield and a pendulum safety toy \\
24 & Capstone: safe candidate-matching policy & End-to-end: offline-pretrained $\to$ constrained online $\to$ audited \\
25 & Applied synthetic ATS RL enhancement audit & Synthetic ATS logging, OPE, constraints, and audit \\
26 & Causal discovery across biological environments; task-aware data enhancement & Synthetic patient effects, invariance diagnostics, and held-out-environment enhancement ablations \\
\bottomrule
\end{longtable}

\section{Prerequisites}
Fluent PyTorch, comfortable with probability, linear algebra, basic optimization theory.
You do \emph{not} need prior RL experience.

\section{Causal-biology extension: one-week plan}
Week 26 takes about ten hours across five study sessions. Start after week 11;
week 14 supplies additional context about support and extrapolation. The later
policy and LLM tracks are not required. Read its teaching guide, calculate a
confounded and adjusted effect, compare observationally equivalent graphs,
study which mechanisms can change across environments, and evaluate enhancement
with patient-disjoint training, validation, and test sets. The workbook and
mini-project use invented data with known mechanisms.

The extension is motivated by Synlico's scientific discussion on causal
discovery across environments and task-aware enhancement. Primary sources and
their limitations are recorded in \texttt{research/causal\_biology.md}.
Multi-environment S-FCI and invariant causal prediction are reading topics;
the numerical diagnostic in the project is not a reproduction of either
algorithm. Company descriptions of proprietary enhancement do not supply
independent evidence of clinical or target-discovery performance.

\chapter{Bibliography pointers}
The per-week \texttt{extended\_reading.pdf} carries the focused paper list for that lecture.
A consolidated bibliography is maintained at \texttt{research/}.

\end{document}
